% Source: https://artofproblemsolving.com/community/c3h1782319_14dhaka04
% Stable identity: aops-topic-1782319-q01
We need $p^9-8p^3\ne0$, i.e.\ $p^6\ne8$, which holds for all positive integers. Simplify both sides. The left fraction is
\[
\dfrac{p^2(p^2+1)^2}{p^3(p^6-8)}=\dfrac{(p^2+1)^2}{p(p^6-8)},
\]
and $\dfrac{5}{6q}+\dfrac{1}{4q}=\dfrac{10+3}{12q}=\dfrac{13}{12q}$. So the equation is
\[
12q\,(p^2+1)^2=13p\,(p^6-8).
\]
For $p=1$ the right side is negative, so $p\ge2$. Now look at the factor $p^2+1$:
\begin{itemize}
\item $\gcd(p^2+1,p)=1$;
\item $p^2\equiv-1\pmod{p^2+1}$ gives $p^6\equiv-1$, so $p^6-8\equiv-9\pmod{p^2+1}$ and $\gcd(p^2+1,p^6-8)$ divides $9$; but $p^2+1\equiv1$ or $2\pmod3$ is never divisible by $3$, so $\gcd(p^2+1,p^6-8)=1$.
\end{itemize}
Hence $(p^2+1)^2$ is coprime to $p(p^6-8)$, and since it divides $13p(p^6-8)$, it divides $13$. But $(p^2+1)^2\ge25>13$ for $p\ge2$. This contradiction shows that there are \boxed{\text{no positive integers }p,q} satisfying the equation. $\blacksquare$
